% ========================================================
% EXPLICACIÓN DEL ALGORITMO CONCURRENTE Y LA SINCRONIZACIÓN
% ========================================================
\section{Mecanismos de sincronización: el worker pool}\label{sec:sync}

\subsection{Por qué worker pool y no pipeline}

El enunciado ofrece dos patrones: \emph{worker pool} y \emph{pipeline}. Un pipeline divide el
trabajo en \emph{etapas} distintas que se encadenan por canales, y rinde cuando cada dato pasa por
varias transformaciones independientes. Lloyd no tiene esa forma: cada iteración hace una sola cosa
con todos los datos (asignar y acumular) y luego necesita el resultado de \emph{todos} para seguir
(actualizar centroides). Es paralelismo de datos con una barrera, y eso es exactamente un worker
pool: $P$ goroutines iguales que consumen unidades de trabajo homogéneas de un canal, y un punto de
encuentro al final de cada ronda. La Figura~\ref{fig:workerpool-pc2} lo muestra.

\begin{figure}[H]
\caption{Worker pool con reparto por canal y acumuladores por bloque}\label{fig:workerpool-pc2}
\centering
\small
% Figura: worker pool con reparto por canal y acumuladores por bloque (sección 9).
\begin{tikzpicture}[
    font=\footnotesize,
    caja/.style={draw, rounded corners=1pt, align=center, inner sep=3pt, minimum height=9mm},
    canal/.style={draw, rounded corners=1pt, align=center, inner sep=2pt, minimum height=7mm, fill=black!5},
    flecha/.style={-Stealth, semithick}]
  \node[caja, text width=20mm] (main) at (0, 0) {Coordinador\\encola $C$ bloques\\por iteración};
  \node[canal, text width=30mm] (chan) at (3.6, 0) {\texttt{chan int}\\$b_0\;b_1\;b_2\;\cdots\;b_{C-1}$};
  \node[caja, text width=36mm] (w1) at (8.0, 1.35) {Goroutine 1\\recibe $b$, asigna,\\escribe $s_{b}, n_{b}, J_{b}$};
  \node (dots) at (8.0, 0.05) {$\vdots$};
  \node[caja, text width=36mm] (wP) at (8.0, -1.35) {Goroutine $P$\\recibe $b'$, asigna,\\escribe $s_{b'}, n_{b'}, J_{b'}$};
  \node[draw, fill=black!15, minimum width=2.5mm, minimum height=36mm, inner sep=0] (barrera) at (11.3, 0) {};
  \node[above=1mm of barrera, align=center] {\texttt{sync.WaitGroup}\\Add($C$) / Done / Wait};
  \node[caja, text width=26mm] (reduce) at (13.9, 0) {Reducción\\$\sum_{b=0}^{C-1}$ en orden\\$\mu_j = s_j / n_j$};
  \draw[flecha] (main.east) -- (chan.west);
  \draw[flecha] (chan.east) -- (w1.west);
  \draw[flecha] (chan.east) -- (wP.west);
  \draw[flecha] (w1.east) -- (w1.east -| barrera.west);
  \draw[flecha] (wP.east) -- (wP.east -| barrera.west);
  \draw[flecha] (barrera.east) -- (reduce.west);
  \draw[flecha, dashed] (reduce.south) -- ++(0, -2.3)
    -| node[pos=0.25, below] {centroides nuevos: de solo lectura en la siguiente iteración} (main.south);
\end{tikzpicture}

\end{figure}

\subsection{Los tres mecanismos}

\paragraph{Canal de bloques.} \texttt{trabajos} es un \texttt{chan int} con capacidad $C$, la
cantidad de bloques. En cada iteración el coordinador encola los índices $0, \dots, C-1$; las
goroutines los reciben con \texttt{for c := range trabajos} y procesan el rango
$[c \cdot \textit{chunk},\ (c+1) \cdot \textit{chunk})$. El canal hace dos cosas a la vez:
reparte el trabajo y equilibra la carga, porque el worker que termina antes simplemente recibe el
siguiente bloque. Con muchos más bloques que workers, ninguno espera ocioso mientras otro termina
(Sección~\ref{sec:chunk}).

\paragraph{Acumuladores privados por bloque.} Cada bloque $b$ tiene su propio
\texttt{parcial}: sumas $s_{bj}$ ($k \times 6$ flotantes), conteos $n_{bj}$ e inercia $J_b$. Lo
escribe un solo worker por iteración, el que recibió $b$, así que no hay dos goroutines escribiendo
la misma memoria. Las asignaciones $a(i)$ también se escriben sin conflicto: cada bloque cubre un
rango disjunto de $i$.

\paragraph{Barrera con \texttt{sync.WaitGroup}.} Antes de encolar, el coordinador hace
\texttt{barrera.Add(C)}; cada worker hace \texttt{barrera.Done()} al terminar un bloque; el
coordinador hace \texttt{barrera.Wait()}. El \texttt{Add} va \emph{antes} de encolar, no dentro del
worker: si el worker hiciera su propio \texttt{Add}, el \texttt{Wait} podría ejecutarse con el
contador en cero antes de que ningún worker haya empezado, y la barrera dejaría pasar. Es el mismo
error que el curso encontró en el ejemplo de la semana 2, donde un \texttt{WaitGroup} sin
\texttt{Wait} no esperaba a nadie. Un segundo \texttt{WaitGroup}, \texttt{pool}, solo sirve para el
cierre ordenado: al terminar, el coordinador cierra el canal y espera a que las $P$ goroutines
salgan del \texttt{range}.

\subsection{Por qué no hace falta \texttt{sync.Mutex}}

El bucle caliente, asignar y acumular, no comparte estado mutable: los centroides son de solo lectura
durante la asignación, y cada worker escribe solo el parcial de su bloque y las asignaciones de su
rango. El único estado que se escribe de forma compartida son los centroides, y los escribe
únicamente el coordinador, entre el \texttt{Wait} de una iteración y el encolado de la siguiente. Que
eso sea seguro no es una impresión: lo garantiza el modelo de memoria de Go. El \texttt{Done} de cada
worker ocurre antes de que el \texttt{Wait} del coordinador retorne, así que todas las escrituras de
los parciales son visibles cuando empieza la reducción; y el envío de un bloque por el canal ocurre
antes de su recepción, así que cada worker ve los centroides nuevos que el coordinador escribió antes
de encolar. Un mutex protegería un acceso concurrente que, por construcción, no existe. Es el
principio que \textcite{hess2023} y \textcite{xu2021} aplican entre nodos y que la PC1 recogió como
decisión de diseño: sincronizar una vez por ronda, no una vez por dato.

\subsection{Por qué la reducción va en orden de bloque}\label{sec:reproducibilidad}

Sumar en punto flotante no es asociativo: $(a + b) + c$ y $a + (b + c)$ difieren en los últimos
dígitos. Si cada worker acumulara todo lo que le tocó, el resultado dependería de \emph{qué} bloques
le tocaron, o sea del planificador, y dos corridas idénticas darían números distintos. Acumulando por
bloque y reduciendo en orden $b = 0, \dots, C-1$, la versión concurrente da el mismo resultado en
toda corrida con el mismo tamaño de bloque, independiente de $P$ y del planificador.

% Generado por tp/scripts/tablas_informe.py desde tp/reports/*.json. No editar a mano.
\begin{table}[H]
\caption{Inercia final sobre el dataset completo, $k=8$, 10 iteraciones}\label{tab:inercias}
\small
\begin{tabular}{@{}llr@{}}
\toprule
\textbf{Versión} & \textbf{$P$} & \textbf{Inercia final} \\
\midrule
secuencial & --- & 4531798,747970240 \\
concurrente & 1, 2, 4, 8, 16 & 4531798,747970126 \\
\bottomrule
\end{tabular}

{\footnotesize\textit{Nota.} La versión concurrente da el mismo valor para toda cantidad de workers porque reduce los parciales en orden de chunk. La diferencia con la secuencial es el redondeo de sumar en otro orden: $2{,}53 \cdot 10^{-14}$ en términos relativos.\par}
\end{table}


La Tabla~\ref{tab:inercias} lo confirma sobre el dataset completo: la inercia final concurrente es
la misma para $P = 1, 2, 4, 8$ y $16$, en la laptop y en el servidor. Frente a la secuencial la
diferencia es de \cifra{inercia-rel} en términos relativos, el redondeo de sumar en otro orden, que
es exactamente lo que la PC1 anticipó siguiendo a \textcite{he2022} y la razón por la que la
equivalencia se comprueba con tolerancia y no bit a bit. El costo de esta decisión es memoria para
$C$ parciales de $k \cdot 7$ valores; con el bloque por defecto son 173 parciales sobre el dataset
completo, unos 66~KB, frente a los 136~MB de datos.

\subsection{Lo que el diseño no hace}

Se descartaron a propósito, y se mencionan para que la comparación mida concurrencia y no otra
cosa: la poda por desigualdad triangular de Elkan y Hamerly, que \textcite{li2023} usa y que
cambia el algoritmo; \emph{mini-batch}, que cambia el trabajo y la aproximación; y la disposición
por columnas (\emph{SoA}) del arreglo, cuya ventaja depende del compilador. Todas son mejoras
posibles para después, sobre una línea base ya medida.
